\documentclass{article}
\usepackage[UTF8]{ctex}
\usepackage{amsmath,amssymb,amsfonts}
\usepackage{enumitem}
\usepackage{geometry}
\geometry{a4paper, margin=1in}
% === 错题本排版优化（不参与题目标识） ===
\setlength{\parindent}{0pt}
\setlength{\parskip}{0.35em plus 0.10em minus 0.08em}
\setlength{\emergencystretch}{3em}
\setlist{itemsep=0.25em, topsep=0.40em, parsep=0pt}
\allowdisplaybreaks[2]
% === 错题本排版优化结束 ===
\begin{document}
\section*{陕西师范大学2025年数学分析真题}
\section*{188．陕西师范大学2025年数学分析试卷}

\subsection{题目1}
1 ．求极限 \(\lim _{n \rightarrow \infty} \frac{1!+2!+\cdots+n!}{n!}\) 。

\subsection{题目2}
2．求极限 \(\lim _{x \rightarrow+\infty} x^{2}[\ln \arctan (x+1)-\ln \arctan x]\) ．

\subsection{题目3}
3．设 \(f(x)=\int_{0}^{x} \frac{\sin t}{\pi-t} \mathrm{~d} t\) ，求积分 \(\int_{0}^{\pi} f(x) \mathrm{d} x\) ．

\subsection{题目4}
4．计算 \(\iint_{S} \frac{a x \mathrm{~d} y \mathrm{~d} z+(a+z)^{2} \mathrm{~d} x \mathrm{~d} y}{\sqrt{x^{2}+y^{2}+z^{2}}}(a>0), S\) 取半球面 \(z=-\sqrt{a^{2}-x^{2}-y^{2}}\) 的上侧．

\subsection{题目5}
5．计算 \(I=\oint_{L} \frac{x \mathrm{~d} y-y \mathrm{~d} x}{3 x^{2}+4 y^{2}}\) ，其中 \(L\) 是椭圆 \(2 x^{2}+3 y^{2}=1\) ，方向沿逆时针方向。

\subsection{题目6}
6．求数项级数 \(\sum_{n=1}^{\infty} \frac{(-1)^{n+1}}{n(n+1)}\) 的和．

\subsection{题目7}
7．设 \(f(x, y)=\sqrt[3]{x^{3}+y^{3}}\) ．

\begin{itemize}
  \item[（1）] 讨论 \(f(x, y)\) 在原点处沿任何方向的导数是否均存在？
  \item[（2）] 讨论 \(f_{x}(0,0), f_{y}(0,0)\) 是否存在？若存在，求出其值．
  \item[（3）] 讨论 \(f(x, y)\) 在 \((0,0)\) 处是否可微？
\end{itemize}

\subsection{题目8}
8．求方程 \(2 x^{2}+2 y^{2}+z^{2}+8 z y-z+8=0\) 所确定的隐函数 \(z=z(x, y)\) 的极值．

\subsection{题目9}
9．设 \(u_{n} \neq 0\) ，且 \(\lim _{n \rightarrow \infty} \frac{n}{u_{n}}=1\) 。讨论级数 \(\sum_{n=1}^{\infty}(-1)^{n+1}\left(\frac{1}{u_{n}}+\frac{1}{u_{n+1}}\right)\) 是否绝对收敛。

\subsection{题目10}
10．设 \(f(x)\) 在 \([a,+\infty)\) 上连续，当 \(x \rightarrow+\infty\) 时，\(y=2 x\) 是 \(f(x)\) 的渐近线．证明：\(f(x)\) 在 \([a,+\infty)\)上一致连续。

\subsection{题目11}
11．设函数 \(f(x)\) 在 \([0,1]\) 上二阶可导，且 \(f^{\prime \prime}(x) \geq 0\) ．证明： \(\int_{0}^{1} f\left(x^{0}\right) \mathrm{d} x \geq f\left(\frac{1}{7}\right)\) ．

\subsection{题目12}
12．设 \(f(x)\) 在 \([0,+\infty)\) 上单调减，且 \(\lim _{x \rightarrow+\infty} f(x)=0\) ．证明：若 \(f^{\prime}(x)\) 在 \([0,+\infty)\) 上连续，则反常积分 \(\int_{0}^{+\infty} f^{\prime}(x) \sin ^{2} x \mathrm{~d} x\) 收敛。

\subsection{题目13}
13．设 \(f(x)=\sum_{n=1}^{\infty} \frac{\cos n x}{\sqrt{n^{3}+n}}\) ．

\begin{itemize}
  \item[（1）] 证明：\(f(x)\) 在 \((-\infty,+\infty)\) 上连续．
  \item[（2）] 证明：\(F\left(\frac{\pi}{2}\right)<\frac{\sqrt{2}}{2}\) ，其中 \(F(x)=\int_{0}^{x} f(t) \mathrm{d} t\) ．
\end{itemize}

\subsection{题目14}
14．证明 \(I(y)=\int_{0}^{+\infty} e^{-p x} \cos x y \mathrm{~d} x\) 在 \((-\infty,+\infty)\) 上一致收敛，其中 \(p>0\) 。

\subsection{题目15}
15．解答如下问题：

\begin{itemize}
  \item[（1）] 计算积分 \(A=\int_{0}^{1} \int_{0}^{1}\left|x y-\frac{1}{4}\right| \mathrm{d} x \mathrm{~d} y\) ．
  \item[（2）] 设 \(z=f(x, y)\) 在区域 \(D:[0,1] \times[0,1]\) 上连续，且满足条件
\[
\iint_{D} f(x, y) \mathrm{d} x \mathrm{~d} y=0 \text { 和 } \iint_{D} x y f(x, y) \mathrm{d} x \mathrm{~d} y=1 \text {. }
\]
证明：存在 \((\xi, \eta) \in D\) ，使得 \(|f(\xi, \eta)| \geq \frac{1}{A}\) ，其中 \(A\) 为（1）中的积分值．
\end{itemize}
\end{document}
